% =========================================================
% Proof: Complement Characterization of Nowhere-Dense
% Source: notes/notes-nowhere-dense-baire.tex
% Difficulty: Easy
% =========================================================

\subsection*{Complement Characterization of Nowhere-Dense}
\label{prf:nowhere-dense-equiv}

\begin{remark*}[Return]
$\leftarrow$ \hyperref[prop:nowhere-dense-equiv]{Back to Proposition in Notes}
\end{remark*}

\begin{proof}
\Claim $E \subseteq X$ is nowhere-dense if and only if
$X \setminus \overline{E}$ is dense in $X$.

\Given A metric space $(X,d)$ and a subset $E \subseteq X$.

\Goal To show $\operatorname{int}(\overline{E}) = \varnothing
\iff \overline{X \setminus \overline{E}} = X$.

\Strategy Unpack both sides in terms of open sets.
$\operatorname{int}(\overline{E}) = \varnothing$ means no non-empty open set
is contained in $\overline{E}$, i.e.\ every non-empty open set meets $X \setminus \overline{E}$.
By Proposition~\ref{prop:dense-equiv}(iii), this is exactly density of $X \setminus \overline{E}$.

% -------------------------------------------------------
% YOUR PROOF HERE
% -------------------------------------------------------

\end{proof}

\begin{remark*}[Proof shape]
Both conditions say the same thing about open sets: translate
``empty interior of $\overline{E}$'' into ``every open set meets the complement,''
which is the density criterion.

\FlashProof{Both conditions say the same thing about open sets: translate ``empty interior of $\overline{E}$'' into ``every open set meets the complement,'' which is the density criterion.}
\FlashUses{\begin{itemize} \item Definition~\ref{def:nowhere-dense}: $\operatorname{int}(\overline{E}) = \varnothing$. \item Proposition~\ref{prop:dense-equiv}(iii): dense iff meets every open set. \item Definition of interior: largest open subset of $\overline{E}$. \end{itemize}}
\end{remark*}

\begin{remark*}[Dependencies]
\begin{itemize}
  \item Definition~\ref{def:nowhere-dense}: $\operatorname{int}(\overline{E}) = \varnothing$.
  \item Proposition~\ref{prop:dense-equiv}(iii): dense iff meets every open set.
  \item Definition of interior: largest open subset of $\overline{E}$.
\end{itemize}
\end{remark*}
